% !TEX root = ../main.tex

\chapter*{前言}
\addcontentsline{toc}{chapter}{前言}

\section*{何以「仿罗什风」译阿含？}

《杂阿含经》与巴利《相应部》同为早期结集之重要见证，但求那跋陀罗译（大正第 99 册）往往直译、冗复，
偶有误读，读来多碍。鸠摩罗什所译大乘经典以删繁就简、义译圆通著称——若\textbf{同一译场精神}
（非同一历史事实）面对\textbf{阿含之教}，文体当如何？

本书即此一\textbf{推想性的文体实验}：义理以巴利平行经为主据（信），
古文主文仿罗什节奏（雅），并附今译意（达）。成书不冒充罗什译经，而供学界与读者\textbf{对照检验}。

\section*{与现有文献之关系}

\begin{itemize}
  \item \textbf{大正《杂阿含经》}：底本与经号锚点，非义理之最终标准；亦非本册之通读顺序。
  \item \textbf{巴利《相应部》《增支部》等}：厘定法义之主要依据。
  \item \textbf{比较研究}（如印顺法师、阿纳罗、宾根海默等）：本册经序重排与之可互参；
    本书侧重\textbf{全经系统化译注}与\textbf{按重排卷次之对照阅读}。
  \item \textbf{本系列 V1}：同一译文之《研究译注本（大正卷序）》，按传统经号与五十卷编排。
\end{itemize}

\section*{关于经序重排}

求那跋陀罗译《杂阿含经》流传久远，\textbf{经号、品目（相应）、卷次}之间时有乖违，
学界通称\textbf{错简}：底本偶见\textbf{错置的卷题}，或\textbf{非本集经叙}被\textbf{插入}经序之中。
Anesaki（1908）及印顺《杂阿含经论会编》等提出依相应结构之\textbf{卷次重排}假说，
便于按早期结集逻辑阅读。

\textbf{本册（V2）的处理原则：}
\begin{itemize}
  \item \textbf{重排阅读序}：依 Anesaki／印顺所采用之 48 个大正卷次顺序编排「学术卷」；
    每经保留大正经号，并标学术序。
  \item \textbf{插入入附录}：第 604、640--641 经（阿育王传说）不入正编阅读序。
  \item \textbf{卷二十三、二十五}：学界多视为阿育王传替补卷，不在上述 48 卷表中；
    本阶段将其中非插入经文暂置于正编末（保证不缺号），后续对照会编再微调。
  \item \textbf{底本}对照栏仍剥除错置卷题；插入经题行保留
    \textit{〔T99插入·非相应经〕} 标注。
  \item \textbf{不}声称此序即唯一「原典」；读者若需馆藏编号，请用大正经号或改读 V1。
\end{itemize}

\noindent 详表与政策见项目文档 \texttt{docs/V2\_ORDER.md}。

\section*{成书状态与局限}

\begin{itemize}
  \item 1362 经已完成首译；216 经为底本略式之重建，须与 1146 经完整译本分档阅读。
  \item 主文与今译意由大语言模型在平行约束下生成并经审读，详见「出版说明·制作方式」。
  \item 末段部分经无巴利专经平行，校勘信度多为中等——依汉本与早期教理常识，宜保守看待。
  \item 印顺 51 相应与 SA 区间之细表仍在完善中；本册阶段一以卷次重排为主。
  \item 本书为\textbf{研究读本}，非终审定本；欢迎依脚注与校勘说明勘误。
\end{itemize}

\section*{致读者}

若您使用本书，建议：先读【正文】与【校勘】，再核【底本】与【平行】；
引用时请标\textbf{大正经号}（可兼标学术序）；引用 ◇ 所标重建经时，请注明「据平行经重建」。
电子版对照阅读器可切换 V1／V2；进一步文献说明见「文献与制作说明」。

\vfill
\begin{flushright}
《杂阿含经·仿罗什风新译》项目组\\
2026
\end{flushright}

\clearpage
